\documentclass[11pt,a4paper]{article}
\usepackage{fontspec,amsmath,amssymb,amsthm,microtype}
\setmainfont{lmroman10-regular.otf}[BoldFont=lmroman10-bold.otf,ItalicFont=lmroman10-italic.otf,BoldItalicFont=lmroman10-bolditalic.otf]
\usepackage[margin=25mm]{geometry}
\usepackage[unicode,hidelinks]{hyperref}
\hypersetup{pdfauthor={LangevinLean Project}}
\setlength{\emergencystretch}{3em}
\setlength{\parskip}{0.4em}
\setlength{\parindent}{0pt}
\allowdisplaybreaks[2]
\raggedbottom
\title{\texorpdfstring{Weak Gaussian Taylor expansion with exactly two derivatives}{Weak Gaussian Taylor expansion with exactly two derivatives}}
\author{LangevinLean Project}
\date{}
\begin{document}\maketitle
Draft mathematical exposition; these notes may contain mistakes.


The potential in the Laplacian theorem (Theorem 2, source th:3) is \(C^{3}\), so its gradient is only \(C^{2}\). These estimates establish the necessary compact-test weak expansion with precisely \(C^{2}\) regularity. A uniformly continuous compact Hessian supplies a freely small quadratic remainder plus a cubic bound; a third derivative of the test is never assumed. The Gaussian averaging requires only a second moment of the random drift.


\section{\texorpdfstring{Express the remainder using a genuine Hessian difference}{Express the remainder using a genuine Hessian difference}}


Restrict the test to the line \(x+t\) \(v\). The ordinary Taylor theorem with first-order polynomial and second derivative remainder gives a point \(u\) strictly between zero and one. The first derivative on this line is the Frechet derivative applied to \(v\); the second derivative is the actual second Frechet derivative applied twice to \(v\). Subtract the quadratic Taylor term at \(x\). The remaining error is exactly one half of the difference between the Hessians at \(x+u\) \(v\) and \(x\), evaluated twice on \(v\).



Earlier results used here: \cite[Section~3.55.1]{LangevinLean2026ULASC}.

\section{\texorpdfstring{Use uniform Hessian continuity without a third derivative}{Use uniform Hessian continuity without a third derivative}}


A compactly supported \(C^{2}\) test has a continuous compactly supported Hessian, hence a uniformly continuous and bounded Hessian. Given \(\eta\) positive, choose \(\delta\) so that Hessian differences at distance less than \(\delta\) are below \(2\eta\). If the displacement norm is below \(\delta\), the segment remainder is at most \(\eta\) times its squared norm. For larger displacements, the uniform Hessian bound gives a quadratic error, which is bounded by a fixed cubic multiple because the displacement norm is at least \(\delta\). Combining the cases proves

\[|R_2(x,v)|\le\eta\|v\|^2+C_\eta\|v\|^3.\]

The same bounded Hessian proves that the actual gradient of the compact \(C^{2}\) test is Lipschitz.



Earlier results used here: \cite[Section~3.55.1]{LangevinLean2026ULASC}.

\section{\texorpdfstring{Separate the deterministic drift from Gaussian noise}{Separate the deterministic drift from Gaussian noise}}


Expand the test first at the Gaussian perturbation alone. Its linear term averages to zero and its quadratic term averages to \(h\) times the true Laplacian. Integrating the \(C^{2}\) remainder uses the exact Gaussian second moment and finite Gaussian third moment. Choose \(\eta\) small enough that the quadratic Gaussian contribution is at most \(\varepsilon\) \(h\). Insert the deterministic \(hv\) shift afterward using the first-order Lipschitz-gradient remainder. This contributes only \(h^2\) times the squared drift norm and a mixed \(h^{3/2}\) term. For \(h\) at most one, their sum is bounded by

\[h(\varepsilon+K_\varepsilon\sqrt h)(1+\|v\|^2).\]

The bound is uniform in the center, drift and admissible step size. No third moment of the drift is required.



Earlier results used here: \cite[Section~3.55.1]{LangevinLean2026ULASC}.

\section{\texorpdfstring{Lift the uniform bound through bounded past weights}{Lift the uniform bound through bounded past weights}}


The \(C^{2}\) Laplacian is continuous as a finite sum of continuous second derivatives. This makes the entire weak residual measurable for measurable random centers and drifts. If the past weight has absolute value at most \(A\) almost everywhere, multiply the pointwise residual bound by \(A\) and integrate. Square integrability of the random drift makes the majorant integrable. The resulting bound is \(Ah(\varepsilon+K\sqrt h)\) times one plus the drift second moment. The fixed-\(K\) version retains a supplied uniform pointwise constant unchanged, so the same constant can be used across every Euler mesh and mesh state.


\begin{thebibliography}{ULA-LSI}
\bibitem[ULA-SC]{LangevinLean2026ULASC}
LangevinLean Project. \emph{Unadjusted Langevin under Strong Convexity}. LL-01, 2026. Draft report, content revision 263fb89cdbf0. \url{https://odakyildiz.github.io/LangevinLean/docs/pdf/UnadjustedLangevin/FullProof.pdf}.
\end{thebibliography}

\end{document}
